\documentclass{article}

\usepackage[margin=1in]{geometry}
\usepackage{hyperref}
\usepackage{listings}
\usepackage{graphicx}
\usepackage{booktabs}
\usepackage{enumitem}
\usepackage{xcolor}

\hypersetup{
  colorlinks=true,
  linkcolor=blue,
  urlcolor=blue,
  pdftitle={Level 1 -- Foundation Chatbot (Beginner)}
}

\lstset{
  basicstyle=\ttfamily\small,
  breaklines=true,
  frame=single,
  backgroundcolor=\color{gray!10}
}

\title{Level 1 -- Foundation Chatbot (Beginner)}
\author{Swaroop Formulation Industries}
\date{2026-03-25}

\begin{document}
\maketitle

\section{Goal}

Deliver a \textbf{working chatbot} on a \textbf{static site}, deployed at no cost on \textbf{GitHub Pages} or \textbf{Vercel}.

\section{Tech Stack}

\begin{tabular}{@{}ll@{}}
\toprule
\textbf{Area} & \textbf{Choice} \\
\midrule
Frontend & \textbf{Vite} + \textbf{React} (TypeScript) \\
LLM & \textbf{Groq API} (free tier, Llama 3) \\
Hosting & \textbf{Vercel} free tier \\
\bottomrule
\end{tabular}

\section{Features}

\begin{itemize}[leftmargin=*]
  \item Single-page site with \textbf{Swaroop} branding.
  \item \textbf{Embedded chatbot widget} (bottom-right).
  \item Chatbot answers questions about \textbf{biodegradable bags}, \textbf{PLA}, \textbf{certifications}, and \textbf{company} information.
  \item \textbf{System prompt} seeded with project report content for grounded tone and facts.
\end{itemize}

\section{Project Structure (Example)}

\begin{lstlisting}[language={},basicstyle=\ttfamily\footnotesize]
src/
  components/
    Chatbot.tsx
  pages/
    index.tsx
  utils/
    llmClient.ts
  styles/
    (Tailwind / global CSS)
public/
api/                    # Vercel serverless (if used)
  chat.ts
\end{lstlisting}

Additional files typically include \texttt{vite.config.ts}, \texttt{package.json}, \texttt{tailwind.config.js}, and environment templates.

\section{Implementation Steps}

\begin{enumerate}[leftmargin=*]
  \item \textbf{Initialize} a Vite + React + TypeScript project.
  \item \textbf{Create} the landing page with company information and branding.
  \item \textbf{Build} the chatbot component with local message state and UI scroll behavior.
  \item \textbf{Integrate} the Groq API via a \textbf{serverless function} (proxy) so the API key stays server-side.
  \item \textbf{Style} with \textbf{Tailwind CSS} for layout, widget position, and responsive behavior.
  \item \textbf{Deploy} to Vercel (connect repository, set env vars, verify build).
\end{enumerate}

\section{Deployment}

\subsection{Vercel (recommended)}

\begin{itemize}[leftmargin=*]
  \item \textbf{Free tier} includes generous \textbf{bandwidth} (on the order of \textbf{100 GB/month} for typical hobby use; confirm current limits in Vercel docs).
  \item \textbf{Serverless functions} for API routes that call Groq without exposing keys in the browser.
\end{itemize}

\subsection{GitHub Pages (alternative)}

Host static assets only; you still need a separate backend or serverless endpoint for the LLM unless using a client-only demo with a public proxy (not recommended for production keys).

\section{Environment Variables}

\begin{tabular}{@{}lp{8cm}@{}}
\toprule
\textbf{Variable} & \textbf{Purpose} \\
\midrule
\texttt{GROQ\_API\_KEY} & Authenticate requests to Groq (set in Vercel project settings, never commit). \\
\bottomrule
\end{tabular}

\section{Version Control}

\begin{itemize}[leftmargin=*]
  \item \texttt{git init} and a thorough \textbf{\texttt{.gitignore}} (node modules, \texttt{.env}, build output).
  \item Branches: \textbf{\texttt{main}} (stable) + \textbf{\texttt{dev}} (integration).
  \item Tag release: \textbf{\texttt{v0.1.0-l1-chatbot-basic}}.
\end{itemize}

\section{Hints}

\begin{itemize}[leftmargin=*]
  \item Use \textbf{Vercel serverless functions} to \textbf{proxy} the Groq API and keep \texttt{GROQ\_API\_KEY} off the client.
  \item Keep the \textbf{system prompt} under roughly \textbf{4000 tokens} to leave room for user messages and the model context window.
\end{itemize}

\end{document}
